\scanpage{72}
\figure{\begin{tikzpicture}[every node/.style={text=gray}]
\node[xscale=-1,transform shape]at(0,0){$\unsure{3}\ \unsure{=|-}\ \unsure{L_K}$};
\node at(-.1,-.65){$\unsure{r}$};
\end{tikzpicture}}
\figure{\begin{tikzpicture}
\node[rotate=180,transform shape,inner sep=0pt]{\begin{minipage}{15cm}\ttfamily
\begin{tabular}{ll}
UNIVERSITE DES SCIENCES \& TECHNIQUES&Session de Juin 1978\\
\hspace*{2cm}DU LANGUEDOC&\\
\hspace*{2cm}\uline{Mathématiques}&
\end{tabular}
\par\vspace{1.1cm}
\begin{center}C4 Math. Pures et Approfondies,\\[12pt]
\uline{Option de M. Petrich}\end{center}
\vspace{1.1cm}
\hspace*{1.2cm}Décrire les classes de conjugaison de sous-groupes et d'éléments

du groupe alterné $\unsure{G_5|\mathcal G_5}$, en indiquant leur interprétation géométrique en

termes de l'icosaèdre.
\hfill\raisebox{1.5\baselineskip}{-}
\par\vspace{1cm}
\begin{center}\begin{tikzpicture}\draw(0,0)--(1.3,0);\end{tikzpicture}\end{center}
\end{minipage}};
\end{tikzpicture}}
\mspage{71}
